%
% 2-geometrie.tex
%
% (c) 2025 Prof Dr Andreas Müller
%
\section{Geometrische Eigenschaften holomorpher Funktionen
\label{buch:holomorph:section:geometrie}}
\kopfrechts{Geometrische Eigenschaften holomorpher Funktionen}

\input{chapters/020-holomorph/fig/fig-konform.tex}

%
% Konforme Abbildungen
%
\subsection{Konforme Abbildungen}
Wir betrachten eine holomorphe Abbildung $f\colon U\to\mathbb{C}$ einer
offenen Menge nach $\mathbb{C}$.
Ein Punkt $z_0 \in U$ wird auf $f(z_0)\in\mathbb{C}$ abgebildet.
Die Abbildung~\ref{buch:holomorph:fig:konform} zeigt eine solche Abbildung.
Eine Gerade durch den Punkt $z_0$ kann durch die Richtung $d\in\mathbb{C}$
mit $|d|=1$ geschrieben werden.
Die Gerade besteht aus den Punkten $z_0+dt$ mit $t\in\mathbb{R}$.
Die Bildkurven sind $t\mapsto f(z_0+dt)$.
Die Tangente an die Bildkurve im Punkt $f(z_0)$ hat gemäss der
Ableitung  die Richtung von
\[
f'(z_0) \cdot d.
\]
Da Multiplikation mit $f'(z_0)$ eine Drehstreckung beschreibt, werden
Geraden durch $z_0$, die einen Winkel $\omega$ einschliessen, auf zwei
Tangenten abgebildet, die ebenfalls den Winkel $\omega$ einschliessen.
Folglich bleiben Winkel bei der Abbildung $f$ erhalten.
Solche Abbildungen heissen \emph{konform}.
\index{konform}%

\begin{satz}
Eine holomorphe Funktion $f\colon U\to\mathbb{C}$ beschreibt in jedem
Punkt $z$, in dem $f'(z)\ne 0$ ist, eine konforme Abbildung.
\end{satz}

Die Gitterlinien stehen im Definitionsgebiet senkrecht aufeinander.
Die rechten Winkel werden von $f$ auf rechte Winkel zwischen den
Bildkurven der Gitterlinien abgebildet.
Daher schneiden sich die Gitterlinien in
Abbildung~\ref{buch:holomorph:fig:konform}
orthogonal.

%
% Der Abbildungssatz von Riemann
%
\subsection{Der Abbildungssatz von Riemann}

